% GigaScience Data Note: Desert microbiome and geochemistry dataset
% from the Rub' al-Khali (Empty Quarter), Saudi Arabia.
%
% This paper describes the DATA: what was collected, how, where, and
% how to access it. The semantic modelling and knowledge graph
% engineering are described in a companion paper (SWJ).

\documentclass[11pt,a4paper]{article}
\usepackage[utf8]{inputenc}
\usepackage[T1]{fontenc}
\usepackage[margin=1in]{geometry}
\usepackage{graphicx}
\usepackage{booktabs}
\usepackage{longtable}
\usepackage{hyperref}
\usepackage{gensymb}
\usepackage{amsmath}
\usepackage{algorithm}
\usepackage{algpseudocode}
\usepackage{xcolor}
\usepackage{cite}

\graphicspath{{../}}

\title{A multi-expedition soil microbiome and geochemistry dataset
  from the Rub' al-Khali desert}

\author{\textit{Author list to be completed}}

\date{}

\begin{document}
\maketitle

\begin{abstract}
\noindent
\textbf{Background.} The Rub' al-Khali (Empty Quarter) is the largest
contiguous sand desert on Earth and an understudied extreme environment
whose microbial communities and biogeochemistry remain poorly
characterized. Integrated, cross-domain data resources for such
ecosystems are scarce: microbial sequencing, soil chemistry, and
environmental observations are typically deposited in disconnected
repositories with incompatible schemas, which limits reuse and
cross-domain analysis.
\textbf{Findings.} We present a coordinated multi-expedition dataset
comprising 16S~rRNA gene amplicon sequencing profiles from 2{,}535 soil
samples collected at 70 sampling sites along a 1{,}043\,km transect
across the Rub' al-Khali during five field expeditions
(2023--2025); X-ray fluorescence (XRF) measurements of 91 chemical
analytes from 249 measurement sessions; GPS-tagged environmental
metadata (temperature, pressure, humidity) recorded at every site
visit; and historical monthly and annual climate records from the
Open-Meteo database. Taxonomic classifications from multiple primer
and classifier workflows are harmonized against NCBI Taxonomy and the
Genome Taxonomy Database (GTDB), yielding 1{,}401{,}008 taxon--sample
abundance observations. All data are integrated into a formal
knowledge graph that encodes the full provenance chain from field
sampling to digital data products; the knowledge graph comprises
approximately 37~million asserted RDF triples.
\textbf{Conclusions.} The dataset is published as raw sequencing data
in the European Nucleotide Archive (umbrella project
\texttt{PRJEB104209}, amplicon project \texttt{PRJEB106069}), as
versioned OWL/RDF modules distributed from the project website and the
Zenodo archive, and via a public SPARQL endpoint at
\url{https://rubalkhali.science/sparql}. All processing code is
released under a free software licence. A companion paper describes
the formal knowledge graph and the Semantic Web engineering that
supports cross-domain querying.

\vspace{0.5em}
\noindent\textbf{Keywords.} desert microbiome; metagenomics;
geochemistry; X-ray fluorescence; 16S~rRNA; knowledge graph;
Rub' al-Khali; Saudi Arabia; FAIR data
\end{abstract}

% ===================================================================
\section{Data description}
% ===================================================================

The Rub' al-Khali is the largest contiguous sand desert on Earth,
covering approximately 650{,}000\,km$^2$ across the southern Arabian
Peninsula~\cite{mandaville_1986}. Despite extreme aridity and surface
temperatures that frequently exceed 50\degree C, the desert harbours
microbial communities whose composition, diversity, and ecological
drivers remain poorly
characterized~\cite{makhalanyane2015}. Predicting ecosystem responses
to climate change in such environments and informing conservation and
land-use policy require coordinated, multi-variable datasets that
capture microbial, chemical, and environmental information at matching
spatial and temporal resolutions~\cite{fierer2017}. Current resources
fall short: microbial sequencing data are typically deposited in
sequence archives, soil chemistry data in laboratory spreadsheets, and
environmental observations in climate databases, with no common
identifiers or data model linking them. This fragmentation blocks the
cross-domain analyses that extreme-environment microbial ecology
requires.

This Data Note presents a coordinated dataset that addresses the gap
for the Rub' al-Khali. The dataset is the product of five field
expeditions conducted across seasons between early 2023 and late
2025, covering 70 sampling sites along a 1{,}043\,km transect from the
western escarpment of the desert to the Saudi Arabia--Oman border. At
each site we collected bulk surface, bulk subsurface (``deep''), and
rhizosphere soil samples in three biological replicates, yielding
2{,}535 soil samples in total. We generated 16S~rRNA amplicon
sequencing profiles for all samples on Illumina MiSeq, processed them
with a standardized workflow (nf-core/Ampliseq + DADA2), and
harmonized the resulting taxonomic classifications across primer and
classifier combinations by aligning NCBI Taxonomy, the Genome Taxonomy
Database (GTDB)~\cite{gtdb}, and the iNaturalist taxonomy into a
single reference. For a subset of the sites we generated X-ray
fluorescence measurements of 91 soil analytes in 249 independent
measurement sessions. At every site visit we recorded temperature,
atmospheric pressure, and relative humidity, and we additionally
retrieved historical climate data from the Open-Meteo
database~\cite{openmeteo} to provide 12{,}936 monthly and 582 annual
climate records.

All data are integrated into a formal, semantic knowledge graph that
links physical samples to their digital data products through an
explicit provenance chain. We used the Semanticscience Integrated
Ontology (SIO)~\cite{sio} as the upper-level framework, grounding
chemical entities in the Chemical Entities of Biological Interest
(ChEBI) ontology~\cite{chebi}, environmental contexts in the
Environment Ontology (ENVO)~\cite{envo}, and taxonomic identifiers in
NCBI Taxonomy~\cite{ncbitaxon}. The knowledge graph comprises
approximately 37 million asserted RDF triples served through a public
SPARQL~1.1 endpoint. This Data Note describes the data acquisition,
processing, validation, and access methods; the semantic modelling,
ontology design patterns, materialization strategy, and reasoning
evaluation are described in a companion paper targeting the Semantic
Web community~\cite{swj_companion}.

% ===================================================================
\section{Methods}
% ===================================================================

% -----------------------------------------------
\subsection{Study area and sampling design}
% -----------------------------------------------

We conducted five sampling expeditions across the Rub' al-Khali along
a transect from the western escarpment (19.15\degree N, 45.17\degree E)
to the eastern border of Saudi Arabia with Oman (20.71\degree N,
54.73\degree E), spanning 1{,}043\,km
(Figure~\ref{fig:transect}). The transect encompasses 60 primary
sampling sites at predetermined geographic coordinates, together with
10 sites of special interest (water wells, a hot spring, an oil spill
site), for a total of 70 sampling sites.

\begin{figure}[!htbp]
  \centering
  \includegraphics[width=0.9\textwidth]{transect_altitude.png}
  \caption{Altitude profile along the sampling transect across the
    Rub' al-Khali, from the western escarpment to the eastern border
    with Oman.}
  \label{fig:transect}
\end{figure}

To capture temporal variability, we conducted expeditions across
different seasons: late winter~2023 (Trip~1, 64 sites visited);
summer~2023 (Trip~2, truncated to 8 sites due to extreme
temperatures); winter~2024 (Trip~3, 65 sites visited); summer~2024
(Trip~4, 60 sites visited); and autumn~2025 (Trip~5, 62 sites
visited). Table~\ref{tab:env_summary} summarizes the observational
coverage. At each site visit we recorded temperature, atmospheric
pressure, and relative humidity with handheld field instruments, and
during one expedition we additionally collected XRF measurements to
characterize the geochemical composition of the soils.

\begin{table}[h]
  \centering
  \caption{Summary of expedition coverage and environmental
    measurements collected at each site visit.}
  \label{tab:env_summary}
  \begin{tabular}{lllc}
    \toprule
    Expedition & Season & Sites visited & Environmental measurements \\
    \midrule
    Trip 1 & Late winter 2023 & 1--60, 4 special & Temp., press., hum. \\
    Trip 2 & Summer 2023      & 1--8             & Temp., press. \\
    Trip 3 & Winter 2024      & 1--60, 5 special & Temp., press., hum. \\
    Trip 4 & Summer 2024      & 1--60            & Temp., hum. \\
    Trip 5 & Autumn 2025      & 1--60, 2 special & Temp., press., hum. \\
    \bottomrule
  \end{tabular}
\end{table}

At each site we collected samples from three distinct soil
compartments: bulk surface soil (0\,cm depth), bulk subsurface soil
(5--10\,cm depth, hereafter ``deep''), and subsurface soil near plant
roots (3--10\,cm depth, hereafter ``rhizosphere''). We obtained three
biological replicates per compartment by placing soil into sterile
tubes, either by scooping soil directly into the tube or using sterile
plastic spoons. We preserved samples at $-20$\degree C in the field
using an Engel MR040F portable freezer; at the end of each expedition
we transferred samples to a $-80$\degree C freezer in the laboratory
until processing.

We included positive and negative controls for each sampling
campaign to control for contamination artefacts. Positive controls
consisted of replicate tubes of Zymo Mock Community Standard
(ZymoBIOMICS, cat.~\#D6305). Negative controls consisted of sterilized
soil treated with UV light for four hours. Controls accompanied every
batch through DNA extraction, library preparation, and sequencing.

% -----------------------------------------------
\subsection{DNA extraction and sequencing}
% -----------------------------------------------

We utilized 0.9--1.2\,g of soil for DNA extraction using either the
PowerLyzer PowerSoil Kit (QIAGEN, cat.~\#12855-50) or the DNeasy
PowerSoil Pro Kit (QIAGEN, cat.~\#47014). We homogenized all samples
with the PowerLyzer 24 Homogenizer and completed extraction either
manually or via the automated QIAcube Connect system. We quantified
eluted DNA using the Qubit HS dsDNA kit (Thermo Fisher
Scientific, cat.~\#Q32854).

We amplified bacterial DNA with the Bakt\_341F~\cite{Fprimer} and
806R~\cite{Rprimer} primer set, targeting the V3--V4 region of the
16S~rRNA gene. We performed amplification using the Multiplex PCR kit
(QIAGEN) and sequenced the resulting amplicon libraries on an Illumina
MiSeq platform in a $2\times 250$ paired-end configuration. All
library preparation and sequencing was performed in the principal
investigator's laboratory at KAUST.

% -----------------------------------------------
\subsection{Sequencing data processing}
% -----------------------------------------------

We processed raw amplicon reads with nf-core/Ampliseq
(v2.14.0)~\cite{nfcore,ampliseq} and QIIME2
(v2024.10)~\cite{qiime2}. We applied the DADA2~\cite{dada2} algorithm
for error correction, chimera removal, and denoising, producing
amplicon sequence variants (ASVs) at single-nucleotide
resolution. We identified taxa using the Genome Taxonomy Database
(GTDB)~\cite{gtdb}. We retained ASVs that were present in at least
three samples with a total abundance greater than 50 reads. We
identified and removed contaminant reads using the QIIME2
\textit{decontam} module based on our negative control samples,
applying a prevalence threshold of 0.52.

% -----------------------------------------------
\subsection{Taxonomy harmonization}
% -----------------------------------------------

To unify taxonomic classifications from different reference databases,
we developed an algorithm that maps and aligns the NCBI Taxonomy
($O_{ncbi}$), the Genome Taxonomy Database ($O_{gtdb}$), and the
iNaturalist taxonomy ($O_{inat}$). The alignment does not imply any
biological commitment but integrates three reference taxonomies for
the purposes of cross-expedition queries, resolving nomenclature
divergence (for example, lineage-specific alphanumeric suffixes in
GTDB) and structural rank conflicts into a consolidated taxonomic
module $O_{final}$.

Let $O = (C, R, E)$ represent a taxonomy comprising a set of classes
$C$, a set of taxonomic ranks $R$, and a set of subclass axioms
$E$. We define a lexical normalization function
$f_{norm}: C \to \Sigma^*$ that maps a class $c \in C$ to a canonical
string by converting the primary label to lowercase, removing
rank-specific prefixes (for example \texttt{d\_\_}), and stripping
alphanumeric lineage suffixes. Let $\mathit{rank}: C \to R$ denote a
function that retrieves the normalized taxonomic rank of a class, and
let $\mathit{parent}: C \to C$ denote a function that retrieves the
direct superclass of $c$.

The workflow initializes with a lexical alignment phase that
generates a candidate seed set
\[
S_{lex} = \{\,(c_{src}, c_{tgt}) \in C_{ncbi} \times C_{gtdb} \mid
  f_{norm}(c_{src}) = f_{norm}(c_{tgt})\,\}.
\]
Lexical matching between NCBI and GTDB produced 27{,}487 seed pairs.
We subsequently partitioned $S_{lex}$ through a structural
verification step that evaluates rank compatibility and topological
support. A seed pair $(c_{src}, c_{tgt}) \in S_{lex}$ is classified as
a strong anchor $A_{strong}$ if $\mathit{rank}(c_{src}) =
\mathit{rank}(c_{tgt})$ and $f_{norm}(\mathit{parent}(c_{src})) =
f_{norm}(\mathit{parent}(c_{tgt}))$. Pairs whose ranks are compatible
but whose parent labels diverge are classified as weak anchors
$A_{weak}$ ($|A_{weak}| = 4{,}347$), which signify potential
nomenclature reclassifications.

We resolve weak anchors using a bottom-up inference strategy based on
the biological heuristic that identical species-level designations
refer to the homologous entity regardless of divergent higher-rank
classifications. For each $(c_{src}, c_{tgt}) \in A_{weak}$, we
generate a structural hint
$h = (\mathit{parent}(c_{src}), \mathit{parent}(c_{tgt}))$. We
aggregate the hints and apply an equivalence axiom to higher-level
taxa if the vote count for a structural hint exceeds a predefined
threshold, yielding 1{,}514 parent hints. We then graft the
iNaturalist taxonomy into the merged NCBI+GTDB module by matching
iNaturalist classes through the same normalization function,
appending equivalence axioms for singular matches and disambiguating
multi-match cases by rank annotation. The full alignment procedure is
given in Algorithm~\ref{alg:tax_align}.

\begin{algorithm}
\caption{Taxonomy alignment and integration workflow.}
\label{alg:tax_align}
\begin{algorithmic}[1]
\State \textbf{Input:} $O_{ncbi}$, $O_{gtdb}$, $O_{inat}$
\State \textbf{Output:} $O_{final}$
\State $S_{lex} \gets \emptyset$; $A_{strong} \gets \emptyset$; $A_{weak} \gets \emptyset$
\State \textit{// Phase 1: lexical anchoring}
\For{$(c_{src}, c_{tgt}) \in C_{ncbi} \times C_{gtdb}$}
    \If{$f_{norm}(c_{src}) = f_{norm}(c_{tgt})$}
        \State $S_{lex} \gets S_{lex} \cup \{(c_{src}, c_{tgt})\}$
    \EndIf
\EndFor
\State \textit{// Phase 2: structural verification}
\For{$(c_{src}, c_{tgt}) \in S_{lex}$}
    \If{$f_{norm}(\mathit{parent}(c_{src})) = f_{norm}(\mathit{parent}(c_{tgt}))$}
        \State $A_{strong} \gets A_{strong} \cup \{(c_{src}, c_{tgt})\}$
    \Else
        \State $A_{weak} \gets A_{weak} \cup \{(c_{src}, c_{tgt})\}$
    \EndIf
\EndFor
\State \textit{// Phase 3: bottom-up resolution}
\State $H \gets \textsc{AggregateParentHints}(A_{weak})$
\State $O_{merged} \gets O_{ncbi} \cup O_{gtdb}$
\For{$(c_{src}, c_{tgt}) \in A_{strong} \cup H$}
    \State $O_{merged} \gets O_{merged} \cup
           \{\textsc{EquivalentClasses}(c_{src}, c_{tgt})\}$
\EndFor
\State $O_{merged} \gets \textsc{DisambiguateHomonyms}(O_{merged})$
\State \textit{// Phase 4: iNaturalist grafting}
\State $O_{final} \gets O_{merged} \cup O_{inat}$
\For{$c_{inat} \in C_{inat}$}
    \State $M \gets \{c \in C_{merged} \mid f_{norm}(c) = f_{norm}(c_{inat})\}$
    \If{$|M| = 1$}
        \State $O_{final} \gets O_{final} \cup
               \{\textsc{EquivalentClasses}(c_{inat}, m \in M)\}$
    \ElsIf{$|M| > 1$}
        \State $m^* \gets \textsc{DisambiguateByRank}(M, \mathit{rank}(c_{inat}))$
        \State $O_{final} \gets O_{final} \cup
               \{\textsc{EquivalentClasses}(c_{inat}, m^*)\}$
    \EndIf
\EndFor
\State \Return $O_{final}$
\end{algorithmic}
\end{algorithm}

The alignment process successfully produced 27{,}487 verified anchors
bridging NCBI and GTDB. Following the primary integration, we anchored
the iNaturalist taxonomy ($O_{inat}$) into the merged
NCBI+GTDB module. The merged module contains 5{,}878 unique taxa
observed across the five expeditions, mapped onto the reference
taxonomies while preserving the original classifier assertions so that
downstream consumers can re-use the mapping or substitute their own.

% -----------------------------------------------
\subsection{X-ray fluorescence analysis}
% -----------------------------------------------

We measured the elemental composition of soil samples with a portable
energy-dispersive X-ray fluorescence (ED-XRF) spectrometer. Each
measurement session generated concentration values for 91 chemical
analytes, including major oxides and trace elements, expressed as mass
percentages. The instrument reports a ``Light Elements'' (LE)
pseudo-analyte that captures the summed contribution of elements below
the detection limit of the ED-XRF method (primarily H, Li, Be, B, C,
N, O, F, Na, and Mg); we carry this through as a separate pseudo-class
in the knowledge graph so that it does not bias the per-element
concentration totals. We applied detector-sensitivity corrections to
the identified elements and scaled the corrected concentrations so
that the final composition totals 100\% (\(\sum\text{final
concentrations} + \text{LE} = 100\%\)). We mapped each analyte to the
Chemical Entities of Biological Interest (ChEBI)
ontology~\cite{chebi} and to the PubChem
database~\cite{pubchem}. Of the 93 monitored parameters, 80 analytes
mapped directly onto ChEBI classes and 92 parameters aligned to
PubChem Compound Identifiers. In total, the XRF data comprise 6{,}674
individual measurement values across 249 measurement sessions.

% -----------------------------------------------
\subsection{Environmental and climate data}
% -----------------------------------------------

We retrieved historical climate data from the Open-Meteo
database~\cite{openmeteo} for each sampling site, yielding 12{,}936
monthly and 582 annual records encompassing temperature, total
precipitation, rainfall, and relative humidity. Field researchers
manually classified each site according to its biome (for example
desert biome, hot desert biome) and specific environmental features
(for example sand dune, saline pan, gravel pavement, desert oasis). We
subsequently mapped each classification to a class in the Environment
Ontology (ENVO)~\cite{envo}, so that every site is linked to a
controlled biome IRI and to one or more environmental-feature IRIs.

% -----------------------------------------------
\subsection{Data integration and structural validation}
% -----------------------------------------------

We integrated all data into a formal knowledge graph by generating
Resource Description Framework (RDF) assertions for every sample,
process, measurement, and environmental observation. The knowledge
graph uses the Semanticscience Integrated Ontology (SIO)~\cite{sio}
as the upper-level vocabulary and models processes, qualities, and
quantities as subclasses of the corresponding SIO
classes. Domain-specific classes are defined as subclasses of SIO
entities; chemical entities are grounded in ChEBI, environmental
contexts in ENVO, and taxonomic identifiers in NCBI Taxonomy. The
resulting knowledge graph is released as nine modular OWL/RDF files
(see Section~\ref{records}). The modelling details, formal design
patterns, and reasoning workflow are described in the companion
paper~\cite{swj_companion}.

To ensure structural integrity during ingest, we defined Shape
Expressions (ShEx)~\cite{shex} shapes for every core entity type:
sampling sites, soil samples, environmental measurements, DNA
extracts, amplicon libraries, sequencing datasets, quality control
records, and taxonomy observations. Each shape specifies the
mandatory properties, their cardinality, and admissible datatypes
(for example, that every sampling site must declare a label and at
least one ENVO biome link, and that every DNA extract must carry a
quantitative concentration measurement). The validator is implemented
in the Apache Jena ShEx library (version 4.10.0) and is invoked from
the data-generation workflow: every time a new ABox module is produced
it is checked against its shape, and the workflow aborts if any
instance violates the schema. This mechanism acts as a regression
test so that modifications to the extract--transform--load scripts
cannot silently introduce structural errors. All eight shape-annotated
ABox modules passed validation without constraint violations.

We additionally verified the logical consistency of the knowledge
graph using the ELK reasoner~\cite{elk}, an implementation of a
consequence-based reasoning procedure for the OWL~2~EL profile. ELK
confirmed that the ontology is consistent (zero unsatisfiable
classes) both before and after the OWL-Horst materialization step
performed by Virtuoso during deployment. The details of the reasoning
evaluation, including the materialization benchmarks, are reported in
the companion paper~\cite{swj_companion}.

% ===================================================================
\section{Data records}\label{records}
% ===================================================================

All raw sequencing data generated in this study are deposited at the
European Nucleotide Archive (ENA). Project and run accessions are
listed in Table~\ref{tab:ena}. Each FASTQ dataset in the knowledge
graph is cross-referenced to its ENA accession via the
\texttt{rdfs:seeAlso} predicate.

\begin{table}[h]
  \centering
  \caption{ENA project accessions.}
  \label{tab:ena}
  \begin{tabular}{lll}
    \toprule
    Accession & Scope & Link \\
    \midrule
    \texttt{PRJEB104209} & Umbrella project &
      \url{https://www.ebi.ac.uk/ena/browser/view/PRJEB104209} \\
    \texttt{PRJEB106069} & 16S rRNA amplicons &
      \url{https://www.ebi.ac.uk/ena/browser/view/PRJEB106069} \\
    \bottomrule
  \end{tabular}
\end{table}

We distribute the integrated knowledge graph as nine modular OWL/RDF
files listed in Table~\ref{tab:modules}. Each file is self-contained:
consumers interested in a specific data type (for example only the
sampling sites and their biome annotations, or only the XRF data) can
download the single file they need without pulling the full
graph. All modules are available as OWL/XML from the project download
page (\url{https://rubalkhali.science/data}) and are archived with a
persistent identifier in the Zenodo deposit~\cite{zenodo_rak} and as
a curated GigaDB entry~\cite{gigadb_rak}. The SPARQL 1.1 endpoint at
\url{https://rubalkhali.science/sparql} serves the same data for
live querying, and a web portal at \url{https://rubalkhali.science/}
provides an interactive map and per-site dashboards for users who do
not wish to write SPARQL.

\begin{table}[h]
  \centering
  \caption{Distributed knowledge graph modules.}
  \label{tab:modules}
  \small
  \begin{tabular}{lp{0.58\textwidth}}
    \toprule
    File & Contents \\
    \midrule
    \texttt{rubalkhali.owl} & Base ontology (TBox); 302 domain
      classes grounded in SIO, ChEBI, ENVO, PATO, UO. \\
    \texttt{rubalkhali\_sites.owl} & 70 sampling sites with
      coordinates and ENVO biome annotations. \\
    \texttt{rubalkhali\_samples.owl} & 2{,}535 soil samples with
      compartment type (surface, deep, rhizosphere) and replicate
      identifier. \\
    \texttt{rubalkhali\_measurements.owl} & Environmental measurements
      (temperature, pressure, humidity) together with 12{,}936 monthly
      and 582 annual climate records. \\
    \texttt{rubalkhali\_xrf.owl} & 6{,}674 XRF measurement values
      across 91 analytes from 249 measurement sessions. \\
    \texttt{rubalkhali\_dna.owl} & 2{,}410 DNA extracts with
      concentration measurements, extraction agents and kits, and
      protocol links. \\
    \texttt{rubalkhali\_sra.owl} & 1{,}242 sequencing run records
      linked to ENA accessions via \texttt{rdfs:seeAlso}. \\
    \texttt{rubalkhali\_qc.owl} & Sequencing quality metrics (GC
      content, duplicate rate, sequence counts) for every run. \\
    \texttt{rubalkhali\_taxonomy\_abox.owl} & 1{,}401{,}008
      taxon--sample abundance observations with absolute counts and
      relative proportions. \\
    \bottomrule
  \end{tabular}
\end{table}

% ===================================================================
\section{Technical validation}
% ===================================================================

We validated the dataset using four complementary approaches: control
samples processed alongside every sequencing batch; structural
validation of every RDF module against formal Shape Expressions
schemas; logical consistency checking of the integrated knowledge
graph with a description-logic reasoner; and competency question
evaluation via SPARQL queries that emulate the intended query workload.

\paragraph{Control samples.}
Positive and negative controls were processed through the full
extraction, amplification, and sequencing workflow in every
expedition batch. Zymo Mock Community standards yielded the expected
taxonomic composition within the error bounds reported by the
manufacturer, confirming the amplification and classification
pipeline. UV-sterilized negative controls carried no reproducible
off-target sequences; contaminant ASVs flagged by the QIIME2
\textit{decontam} module (prevalence threshold 0.52) were removed
prior to downstream integration.

\paragraph{Structural validation.}
We integrated the Apache Jena ShEx validator into the automated
data-generation workflow so that every ABox module is tested against
its shape immediately after generation and before loading into the
triple store. The validator checks that mandatory properties are
present, that cardinalities are respected, and that datatypes are
consistent with the expected domain. The workflow aborts with an
explicit error message if any instance violates a shape constraint,
for example if a sampling site lacks GPS coordinates or a DNA extract
lacks a concentration value. All eight shape-annotated ABox modules
(\texttt{rubalkhali\_sites.owl}, \texttt{rubalkhali\_samples.owl},
\texttt{rubalkhali\_measurements.owl}, \texttt{rubalkhali\_dna.owl},
\texttt{rubalkhali\_xrf.owl}, \texttt{rubalkhali\_sra.owl},
\texttt{rubalkhali\_qc.owl}, and \texttt{rubalkhali\_taxonomy\_abox.owl})
passed ShEx validation without constraint violations.

\paragraph{Logical consistency.}
We verified the logical consistency of the full knowledge graph
(base ontology, all ABox modules, and the taxonomy reference module)
with the ELK reasoner~\cite{elk}. ELK confirmed that the ontology is
consistent---no class is reduced to the empty set and no individual
belongs to disjoint classes---both before and after OWL-Horst
forward-chaining materialization has been applied by the deployment
engine. We chose ELK over more expressive reasoners such as
HermiT~\cite{hermit} because the 37-million-triple ABox renders
full OWL~2~DL reasoning computationally infeasible on commodity
hardware. The consequence-based calculus implemented by ELK produces
class-hierarchy results under the OWL~2~EL profile in under two
seconds on the asserted graph, and the post-materialization run
confirms that the additional 5.7~million type assertions do not
introduce disjointness violations. Full reasoning benchmarks are
reported in the companion paper~\cite{swj_companion}.

\paragraph{Competency questions.}
We defined nine competency questions expressed as SPARQL queries with
known expected answers, covering (i)~per-analyte XRF retrieval,
(ii)~site listings with GPS coordinates, (iii)~biome assignments via
OWL restrictions, (iv)~Light Elements percentages, (v)~DNA extraction
records, (vi)~provenance tracing from soil sample through library
preparation and sequencing to ENA run accession, (vii)~taxonomic
composition per sequencing run, (viii)~cross-domain queries joining
environmental variables with taxon abundances, and (ix)~per-sample
quality control metrics. All nine queries return the expected result
sets against the deployed endpoint. The queries are executed as part
of continuous integration so that any regression in the data or in
the ontology surfaces as a test failure; full queries and latency
benchmarks appear in the companion paper~\cite{swj_companion}.

% ===================================================================
\section{Re-use potential}
% ===================================================================

The dataset supports several re-use scenarios. At the smallest scale,
individual investigators can retrieve the taxonomic composition and
accompanying environmental metadata for a specific site or trip by
downloading the relevant ABox modules or by issuing SPARQL queries
against the endpoint. At the largest scale, the full knowledge graph
supports cross-domain queries that would be difficult or impossible to
express against a conventional relational store. For example,
researchers interested in microbial responses to thermal stress can
retrieve the relative abundance of particular taxa at every site
where the recorded field temperature exceeds a threshold, with the
SPARQL query traversing soil sample, DNA extract, amplicon library,
sequencing run, FASTQ dataset, and relative-abundance observation in a
single go. Researchers interested in geochemistry can ask for XRF
concentration patterns grouped by ENVO biome, joining the site's
biome assignment with the per-sample analyte values.

The modular OWL architecture also makes it straightforward to extend
the knowledge graph with new data types. Researchers can introduce
new classes as subclasses of existing SIO entities and generate
corresponding ABox modules; the existing ShEx shapes can be extended
to cover the new classes, and the re-materialization workflow in
Virtuoso rebuilds the inferred triples automatically. Candidate
extensions include metabolomics, remote sensing products, biogeography
observations, and additional seasonal field trips. Because the
identifiers are resolvable and the ontology is grounded in widely-used
upper-level vocabularies, third parties can also link their own data
into the knowledge graph by asserting equivalence or relationship
axioms against our IRIs.

Users who do not wish to write SPARQL can access the data through the
web portal at \url{https://rubalkhali.science/}, which provides an
interactive map of sampling sites and per-site dashboards with
environmental measurements, geochemistry, and taxonomic profiles. The
portal also hosts a SPARQL editor pre-populated with the nine
competency questions as click-to-load examples, together with an
autocomplete-aware query editor (YASGUI)~\cite{yasgui}. The same
editor is used internally as a regression test for the live endpoint.

% ===================================================================
\section{Availability of supporting data and materials}
% ===================================================================

\paragraph{Sequencing reads.}
Raw 16S~rRNA gene amplicon reads are deposited at the European
Nucleotide Archive under the umbrella project \texttt{PRJEB104209}
(\url{https://www.ebi.ac.uk/ena/browser/view/PRJEB104209}) and the
amplicon sub-project \texttt{PRJEB106069}
(\url{https://www.ebi.ac.uk/ena/browser/view/PRJEB106069}).

\paragraph{Knowledge graph modules.}
All OWL/RDF modules listed in Table~\ref{tab:modules}, together with
the ShEx shape definitions, the SPARQL competency questions, the
schema documentation, and the ontology design pattern documentation,
are archived on GigaDB~\cite{gigadb_rak} and on
Zenodo~\cite{zenodo_rak}. The same files are served live from the
project download page at
\url{https://rubalkhali.science/data}. A public SPARQL~1.1 endpoint
is available at \url{https://rubalkhali.science/sparql}.

\paragraph{Source code.}
All source code for data processing, ontology engineering, RDF
generation, ShEx validation, the SPARQL competency questions, the
web portal (React + FastAPI), and the Docker Compose deployment is
publicly available on GitHub at
\url{https://github.com/bio-ontology-research-group/empty-quarter}. The
repository includes Groovy scripts that use the OWL API 5 for
OWL/RDF generation, Python scripts for data processing and
statistical analysis, and the deployment scripts that materialize the
graph in Virtuoso.

\paragraph{License.}
The dataset, knowledge graph, ontology, ShEx shapes, and source code
are released under the Creative Commons Attribution 4.0 International
licence (CC~BY~4.0,
\url{https://creativecommons.org/licenses/by/4.0/}).

% ===================================================================
\section*{Declarations}
% ===================================================================

\paragraph{List of abbreviations.}
ABox, Assertion Box; ASV, amplicon sequence variant; CC~BY, Creative
Commons Attribution; ChEBI, Chemical Entities of Biological Interest;
DL, Description Logic; ED-XRF, energy-dispersive X-ray
fluorescence; ELK, EL reasoner; ENA, European Nucleotide Archive;
ENVO, Environment Ontology; ETL, extract--transform--load; FAIR,
Findable, Accessible, Interoperable, Reusable; GTDB, Genome
Taxonomy Database; IRI, Internationalized Resource Identifier;
LE, Light Elements; NCBI, National Center for Biotechnology
Information; OWL, Web Ontology Language; PATO, Phenotype and Trait
Ontology; QC, quality control; RDF, Resource Description Framework;
SHACL, Shapes Constraint Language; ShEx, Shape Expressions; SIO,
Semanticscience Integrated Ontology; SPARQL, SPARQL Protocol and
RDF Query Language; TBox, Terminological Box; UO, Units Ontology;
XRF, X-ray fluorescence.

\paragraph{Competing interests.}
The authors declare no competing interests.

\paragraph{Funding.}
Funding statement to be completed.

\paragraph{Author contributions.}
Author contributions to be completed per CRediT.

\paragraph{Acknowledgements.}
Acknowledgements to be completed.

% ===================================================================
\bibliographystyle{plain}
\bibliography{../sn-bibliography}

\end{document}
